\section{Active C-RAG}\label{sec:approach}

\subsection{Prior Art}\label{sec:prior_schema}

In the schema of Figure~\ref{fig:system}, CF and C-RAG only use round $0$. CF performs no retrieval, so that $D_0 = \emptyset$, the LLM $G$ writes the claim $y$ from the query $x$ alone, and the subclaim scorer returns the confidence score $s(c \mid x)$ of the model. C-RAG instead retrieves $D_0$ with the query $x$, and its subclaim scorer returns the score $s(c \mid x, D_0)$ in \eqref{eq:crag_score}. In both methods, the decision block keeps every subclaim, i.e., $k_0 = \hat{y}_0$ and $w_0 = z_0 = \emptyset$, the system stops at $t^* = 0$, and the returned set is $\hat{y} = F_{\hat{\tau}}(\hat{y}_0)$.

In practice, the retrieved documents are short passages, such as chunks of an article or single sentences, and we refer to them as passages. The score \eqref{eq:crag_score} has two shortcomings when the answer to a query relies on a bridge entity, that is, an entity that is not referred to in the query. Firstly, the cosine similarity $\mathrm{sim}(c, d)$ measures how close a subclaim and a passage are in topic, not whether the passage supports the subclaim. Secondly, the factor $\mathrm{sim}(x, d)$ penalises passages about the bridge entity, as these are not very similar to the query. To illustrate this, consider the query ``Where was the director of Jaws born?'', for which the LLM produces the subclaims $c_1 =$ ``Jaws was directed by Steven Spielberg'' and $c_2 =$ ``Steven Spielberg was born in Cincinnati, Ohio''. Retrieval with $x$ will likely surface the article on the film Jaws, which supports $c_1$, but not Spielberg's biography, since nothing in the query points to it. As no passage of $D_0$ supports $c_2$, its score is low and $F_{\hat{\tau}}$ removes it, so that the returned set contains the bridge $c_1$ but not $c_2$, which is the subclaim that answers the query. Note that even if Spielberg's biography were retrieved, its low similarity to $x$ would keep the score of $c_2$ low.

\subsection{Base Pipeline}\label{sec:base}

AC-RAG uses round $0$ of C-RAG unchanged, i.e., the initial retrieval, the LLM $G$ and the subclaim extraction, so that with $T = 0$ it differs from C-RAG only in the subclaim scorer. The subclaims of $\hat{y}_0$ are kept in the order $c_1, c_2, \ldots$ in which they appear in $y$, as an answer typically states a bridge subclaim, such as $c_1$, before the subclaims that rely on it.

\subsection{Subclaim Scorer}\label{sec:scorer}

To score whether the evidence supports a subclaim, we replace \eqref{eq:crag_score} by a natural language inference (NLI) model. For a passage $d$ and a subclaim $c$, the NLI model gives the probability $p_{\mathrm{ent}}(d, c)$ that $d$ entails $c$, and the scorer assigns
\begin{equation}
  s_t(c \mid x, D_t) = \max_{d \in D_t} p_{\mathrm{ent}}(d, c).
  \label{eq:nli_score}
\end{equation}
This score is high only if some passage entails $c$, and it does not depend on how the query $x$ is worded. Spielberg's biography, for instance, clearly entails $c_2$, although it is not similar to $x$. Entailment-based scorers have also been found to match or outperform LLM confidence scores for conformal factuality in RAG, at a far lower computational cost \cite{chen2026conformal}.

The scorer also returns the evidence that each round adds for each subclaim. For a subclaim $c$, let $N_t(c) \subseteq D_t$ denote the passages of $D_t$ that have not been scored against $c$ in earlier rounds. The new evidence for $c$ at round $t$ is then the best of these passages,
\begin{equation}
  p_t(c) = \max_{d \in N_t(c)} p_{\mathrm{ent}}(d, c),
  \label{eq:new_score}
\end{equation}
which is only defined when $N_t(c) \neq \emptyset$. A passage therefore contributes to $p_t(c)$ at most once for each subclaim, including a passage that was retrieved for a different subclaim.

\subsection{Decision Block}\label{sec:decision}

Taking the null hypothesis $H_c$ for each subclaim $c$ to be that $c$ is unsupported, i.e., $a(c \mid x) = 0$, accepting $c$ as supported is equivalent to rejecting $H_c$, which we call certifying $c$. The guarantee \eqref{eq:cf_seq} therefore bounds the family-wise error rate (FWER) over the subclaims of a single answer, i.e., the probability of certifying at least one unsupported subclaim. Since the decision block is applied at every round, and the revision creates new subclaims from the evidence of earlier rounds, both the hypotheses and the number of tests depend on the data.

The decision block first removes the subclaims that the evidence contradicts. To this end, the NLI model also predicts, for each pair of a passage $d$ and a subclaim $c$, whether $d$ entails $c$, contradicts $c$, or neither. A subclaim that is not certified is removed when at least one passage contradicts it and no passage entails it, i.e.,
\begin{equation}
  z_t = \bigl\{ c \in \hat{y}_t \setminus k_{t-1} : \text{some } d \in D_t \text{ contradicts } c, \text{ and no } d \in D_t \text{ entails } c \bigr\},
  \label{eq:removal}
\end{equation}
with $k_{-1} = \emptyset$. Removal requires no calibration, as a removed subclaim is never certified, and so cannot lead to an error in \eqref{eq:cf_seq}. It may, however, reduce the usefulness \eqref{eq:usefulness} when a supported subclaim is removed. Removed subclaims are not reintroduced by the revision block.

For the remaining subclaims, the decision block accumulates evidence over the rounds with an e-process \cite{ramdas2023game}. Its validity rests on a bound on how high the new evidence for an unsupported subclaim scores. Specifically, we require a constant $\bar{p}_{\mathrm{null}} \in (0, 1)$ such that, for every subclaim $c$ with $a(c \mid x) = 0$ and every round $t$ at which $p_t(c)$ is defined,
\begin{equation}
  \mathbb{E}\bigl[\, p_t(c) \mid \mathcal{F}_{t-1} \,\bigr] \leq \bar{p}_{\mathrm{null}},
  \label{eq:null_bound}
\end{equation}
where $\mathcal{F}_{t-1}$ denotes everything observed up to round $t-1$. In other words, $\bar{p}_{\mathrm{null}}$ is the level up to which, on average, the NLI model is fooled by the new evidence for an unsupported subclaim. For each subclaim $c$, the e-process $M_t(c)$ starts at $1$ at the round at which $c$ is created, and at every round at which $p_t(c)$ is defined, it is updated as
\begin{equation}
  M_t(c) = M_{t-1}(c) \bigl(1 + \lambda_t(c) \bigl(p_t(c) - \bar{p}_{\mathrm{null}}\bigr)\bigr),
  \label{eq:betting}
\end{equation}
while it is left unchanged at the other rounds. The stake $\lambda_t(c) \in [0, 1/\bar{p}_{\mathrm{null}}]$ is chosen from $\mathcal{F}_{t-1}$, for instance as a constant or adaptively as in \cite{waudbysmith2023betting}, and its upper limit keeps $M_t(c)$ non-negative. Under \eqref{eq:null_bound}, each factor in \eqref{eq:betting} has expectation at most $1$ given $\mathcal{F}_{t-1}$ for an unsupported subclaim, so that $\mathbb{E}[M_t(c) \mid \mathcal{F}_{t-1}] \leq M_{t-1}(c)$, whatever happened in earlier rounds. By Ville's inequality \cite{ramdas2023game}, the probability that $M_t(c)$ ever reaches $1/\alpha$ is then at most $\alpha$, regardless of the number of rounds and of the round at which the system stops. For a supported subclaim, on the other hand, the new passages typically score above $\bar{p}_{\mathrm{null}}$, and $M_t(c)$ grows multiplicatively over the rounds.

The constant $\bar{p}_{\mathrm{null}}$ is estimated from a pilot set $\mathcal{D}_{\mathrm{pilot}}$ of $m$ labelled queries, which we refer to as pilot queries. The pilot queries are processed by the same loop as any other query, and $\bar{p}_{\mathrm{null}}$ is set to an upper confidence bound on the mean of $p_t(c)$ over their unsupported subclaims, so that \eqref{eq:null_bound} holds with high probability over the pilot set. Note that the pilot set only estimates the average of $p_t(c)$ over the histories of the pilot queries, whereas \eqref{eq:null_bound} must hold given each history. It is important to note that that the pilot set must be kept separate from the $n$ calibration queries $\mathcal{D}_{\mathrm{cal}}$ of the final filtering step. For the guarantees at the final filtering step to hold at test time, $\bar{p}_{\mathrm{null}}$ and thus the decision rule must have been already fixed so that test-time and calibration time queries receive the same treatment by the final filtering step.

The decision block certifies the subclaims at the full level $\alpha$, following fixed-sequence testing (FST) \cite{bauer1991multiple}, in which hypotheses are tested in an order fixed in advance and a hypothesis is only rejected once all hypotheses before it have been rejected. Accordingly, a subclaim is only certified once every subclaim before it in the order $c_1, c_2, \ldots$ has been certified. Let $j_t$ be the position of the first subclaim whose e-process has not yet reached $1/\alpha$ by round $t$, where a subclaim that has been removed or rewritten counts as not certified. The certified subclaims are then
\begin{equation}
  k_t = \bigl\{ c_j : j < j_t \bigr\}, \qquad j_t = \min\Bigl\{ j : \max_{s \leq t} M_s(c_j) < 1/\alpha, \text{ or } c_j \text{ has been removed or rewritten} \Bigr\},
  \label{eq:certify}
\end{equation}
and every subclaim that is neither certified nor removed is reworked, i.e., $w_t = \hat{y}_t \setminus (k_t \cup z_t)$. As Ville's inequality bounds the probability that $M_t(c)$ ever reaches $1/\alpha$, a certified subclaim stays certified even if its e-process decreases later, which agrees with the subclaims of $k_t$ being kept unchanged by the revision block. All subclaims of $w_t$ accumulate evidence in parallel, so that a subclaim that waits behind an earlier subclaim may already exceed $1/\alpha$ when that subclaim is certified. Note that a single factor in \eqref{eq:betting} is at most $1/\bar{p}_{\mathrm{null}}$, so that, unless $\bar{p}_{\mathrm{null}} \leq \alpha$, a subclaim is only certified after several rounds in which strongly entailing passages are found for it. In the example, $c_1$ is entailed by the Jaws article, so that $M_0(c_1)$ increases and further passages about the film are retrieved for it, whereas no passage of $D_0$ entails $c_2$, so that $M_0(c_2) = 1$ and $c_2$ is reworked.

The decision block stops at the round
\begin{equation}
  t^* = \min\bigl\{ t : w_t = \emptyset, \ \text{or} \ D_{t+1} = D_t, \ \text{or} \ t = T \bigr\},
  \label{eq:stop}
\end{equation}
that is, when every subclaim has been certified or removed, when the rework returns no new passages, so that no further evidence can be gained, or when the budget of $T$ rounds is used up. By Ville's inequality, any stopping rule preserves the validity of the e-processes, and \eqref{eq:stop} only determines how much evidence is collected.

\subsection{Query Generation}\label{sec:query}

For each subclaim $c \in w_t$, an LLM writes the retrieval query $q_t(c)$ from the query $x$, the subclaim $c$ and the evidence $D_t$, so that the query can make use of a bridge entity identified in earlier rounds, and the returned passages are added to the evidence store as in \eqref{eq:store}. In the example, $c_2$ yields a query such as ``Steven Spielberg birthplace'', which retrieves his biography. As a simpler alternative, the subclaim $c$ itself can be used as the retrieval query. In multi-hop C-RAG, the queries are written from the query $x$ and the evidence found so far, so that the retrieval has no reason to look for the evidence of a particular subclaim, whereas the queries of AC-RAG are written from the subclaims that are not yet certified.

\subsection{Subclaim Revision}\label{sec:revision}

The revision block keeps the certified subclaims $k_t$ unchanged, and, for each subclaim $c \in w_t$, either leaves $c$ unchanged or rewrites it given the new evidence. It may also add new subclaims that the new evidence supports. A rewritten or added subclaim is a new hypothesis, and it is placed after the existing subclaims in the order, with an e-process that starts at $1$. A refinement of a certified subclaim, such as the month of an event whose year is already certified, is also added as a new subclaim, so that the certified subclaim keeps its status. In the example, $c_2$ needs no rewriting, as Spielberg's biography entails it as it stands, so that $c_2$ keeps its e-process, which grows with the new evidence and allows $c_2$ to be certified once $c_1$ is.

Under FST, a subclaim that is rewritten counts as not certified at its position, by \eqref{eq:certify}. Revisions therefore do not help the subclaims after this position to be certified by the e-processes, but they can still increase the number of supported subclaims that pass the final filtering step.

\subsection{Final Filtering Step}\label{sec:final}

The returned set can be obtained in two ways, which we state side by side.

The first is to return the certified subclaims, $\hat{y} = k_{t^*}$. Let $c^*$ be the first unsupported subclaim in the order. Any error in \eqref{eq:cf_seq} requires certifying $c^*$, since every later subclaim is only certified after it, and this in turn requires $M_t(c^*) \geq 1/\alpha$ at some round. By Ville's inequality, this has probability at most $\alpha$, which gives \eqref{eq:cf_seq} without dividing $\alpha$ among the subclaims. The bound also holds when the evidence for $c^*$ was collected before $c^*$ came under test. As new subclaims are placed after the existing ones, $c^*$ is determined at the round at which it is created, which is also the round at which its e-process starts. This guarantee relies on \eqref{eq:null_bound}, and so holds with high probability over the pilot set.

The second is to apply the filter \eqref{eq:filter_set} to the remaining subclaims, $\hat{y} = F_{\hat{\tau}}(\hat{y}_{t^*})$, with the score \eqref{eq:nli_score} on $D_{t^*}$, where $\hat{\tau}$ is chosen as in CF, but from calibration queries that have each been processed by the whole AC-RAG loop. As long as the loop is fixed in advance and is the same at calibration and test time, the test query and the calibration queries remain exchangeable, and \eqref{eq:cf_seq} holds regardless of \eqref{eq:null_bound} and of how many subclaims the loop recovers. With $T = 0$, this reduces to C-RAG with the score \eqref{eq:nli_score}.

\subsection{Validity and Assumptions}\label{sec:assumptions}

Ville's inequality makes each e-process valid over any number of rounds, and FST makes the certification valid over the subclaims of an answer at the full level $\alpha$. Both rely on the null bound \eqref{eq:null_bound}, which is the only assumption of the first returned set. It can fail in two ways. The NLI model may be fooled by passages that are closely related to an unsupported subclaim but say something slightly different, which raises the mean of $p_t(c)$ under the null, and the pilot set only estimates this mean on average over the histories of the pilot queries, rather than given each history \tk{Check how often the NLI model is fooled by such passages.}. The second returned set does not rely on \eqref{eq:null_bound}, and only requires that the calibration queries and the test query are exchangeable.

\section{Questions for Discussion}

It would be great to get your input on some of these
\begin{enumerate}[leftmargin=*]
  \item Is the argument in Section~\ref{sec:final} correct that the evidence collected for a subclaim before it comes under test can also be used, provided that its e-process starts when the subclaim is created and new subclaims are placed after the existing ones?
  \item How should subclaims that do not depend on one another be ordered, given that one which cannot be certified early prevents every later subclaim from being certified, even if those are well supported? An alternative to FST is to give each subclaim $c$ a weight $\omega_c$ when it is created, with $\sum_c \omega_c \leq 1$, and to certify $c$ once $M_t(c) \geq 1/(\alpha \omega_c)$. By Ville's inequality and a union bound, this also controls the FWER at level $\alpha$, without any subclaim blocking another, at the cost of a higher threshold for each subclaim.
  \item Should part of $\alpha$ be reserved so that certification can continue after a subclaim is removed or rewritten, as in the online fallback procedure of \cite{tian2021online}?
  \item How should $\bar{p}_{\mathrm{null}}$ be estimated, and is a guarantee that holds with high probability over the pilot set acceptable? Can \eqref{eq:null_bound} be expected to hold given each history, or only on average?
  \item Should the e-process bet once per round on the best new passage, as in \eqref{eq:new_score}, or once per passage? Betting once per passage certifies supported subclaims in fewer rounds, but makes \eqref{eq:null_bound} harder to justify when several passages come from the same article.
  \item Can the certified set $k_{t^*}$ replace the filter $F_{\hat{\tau}}(\hat{y}_{t^*})$, or should both be reported?
\end{enumerate}